\chapterOrPart{Conclusions}\label{ch:conclusions}

\begin{explanation}
    This chapter is mandatory. Your work must have conclusions.
    
    The conclusions contain your main findings and the answer to the research question(s) in a clear and concise way. Conclusions should not contain results that have not yet been discussed in your work.
\end{explanation}

The central answer is that climate forcing and technology-stock assumptions interact asymmetrically in the executed thesis model. Under a frozen technology stock, climate forcing mainly changes the useful thermal side of the residential balance: in the strongest reported case, long-term RCP\,8.5, HDD$_{18}$ falls by about \SI{31.0}{\percent}, CDD$_{22}$ rises by about $94.9\,^{\circ}\mathrm{C}\cdot\mathrm{d}$ per year, and useful space-heating demand falls by about \SI{33.3}{\percent}. This answers the first supporting question: warming reduces heating pressure and raises cooling pressure, but the current model does not convert cooling pressure into active cooling electricity demand. Climate forcing alone therefore changes annual heating indicators much more clearly than it changes peak grid import under the frozen stock; in the cold-design-year comparison, the peak remains essentially unchanged at about \SI{30.6}{kW} versus the historical \SI{30.8}{kW}.

The technology cases dominate the final-energy carrier and peak-load results. In the cold-design-year long-term RCP\,8.5 comparison, frozen stock reduces gas use while leaving peak import close to the historical baseline, but moderate electrification raises cohort peak import to about \SI{106.45}{kW}, and the high-electrification PV--EV case raises it to about \SI{188.40}{kW}. This answers the second supporting question: electrification shifts residential demand from fossil carriers toward electricity and strongly increases grid-import stress when unmanaged, while PV affects self-consumption and export but does not remove the peak problem in the executed high-electrification case. The economic and emissions results are conditional on the configured tariff, CAPEX and fixed emission-factor assumptions: frozen stock lowers the dynamic-tariff bill and operational emissions, moderate electrification lowers operational emissions only slightly while raising bills, and the high-electrification PV--EV case raises both bills and operational emissions under the fixed-factor accounting.

The stochastic and validation evidence set the interpretation boundary for the whole thesis. Household diversity is represented through stochastic cohorts and the Richardson-style validation indicates plausible diversity structure, but the executed design-year results use five stochastic seeds and do not provide a full variance decomposition across climate, technology and behaviour. The third supporting question is therefore only partially answered: technology-driven peak differences are visibly much larger than the frozen-stock climate-only peak effect in the reported stress table, but the thesis should not claim a fully quantified stochastic uncertainty hierarchy. The fourth supporting question is also mixed. Annual calibration converges, PV validation is strong against PVGIS and Elia, and EV timing agrees well with the Fluvius EV signature, but the aggregate Fluvius smart-meter profile comparison fails with CVRMSE around \SI{88}{\percent}, Pearson correlation around $0.40$, and peak-timing error around $18$ hours. The model is consequently defensible as a traceable conditional scenario model for Belgian residential demand, not as an externally calibrated Belgian smart-meter forecasting model.


\mode<all> % do not remove
